\documentclass[../main.tex]{subfiles}

\begin{document}
	\chapter{Fase 5. Desarrollo del Generador de Cartuchos Compactos}
	\label{sec:DerarrolloFase5}
	
	Este anexo recoge el detalle completo del desarrollo de la Fase 5, generando como artefacto cartuchos compactos que se pueden cargar en el ESP32.
	
	\input{capitulos/F05/F5-Fase 5. Formato del Cartucho.tex}
	\input{capitulos/F05/F5-I01. Añalisis}
	\input{capitulos/F05/F5-I02. Formato del Cartucho.tex}
	\input{capitulos/F05/F5-I03. Desarrollo del generador}
	\input{capitulos/F05/F5-I04. Creación del módulo}
	\input{capitulos/F05/F5-I05. Validaciones}	
	\input{capitulos/F05/F5-I06. Creación del GWC}	
	\input{capitulos/F05/F5-I07. Mejoras del manejo}
	\input{capitulos/F05/F5-I08. Las tres imágenes y el GWC.tex}
	\input{capitulos/F05/F5-I09. La lectura unificada de cartuchos}
	\input{capitulos/F05/F5-I10. Pestaña Cartucho}
	\input{capitulos/F05/F5-I11. Manejo de ficheros Empaquetados}
	\input{capitulos/F05/F5-I12. Cambio en el manejo de las pestañas}
	\input{capitulos/F05/F5-I13. El GWC debe ser igual que el GWI}
	\input{capitulos/F05/F5-I14. Cambios menores}
	\input{capitulos/F05/F5-I15. Revisión del motor y del editor}
	\input{capitulos/F05/F5-I15a. Ajustes en la documentación}
	\input{capitulos/F05/F5-I16. Rangos de la UI y mensajes}
	\input{capitulos/F05/F5-I17. Nuevos cambios menores}
	\input{capitulos/F05/F5-I18. Sistema de pruebas}
	\input{capitulos/F05/F5-I19. Carga unificada de cartuchos}
	\input{capitulos/F05/F5-I20. Tamaño de la Memoria}
	\input{capitulos/F05/F5-I21. Datos comunes y su pestaña natural}
	\clearpage
	
\end{document}
